\section*{Lab 6 (Random Forests) Implementation Tasks}
\textbf{Lab 6 focus:} ensemble learning, bagging, feature subsampling, out-of-bag error, feature importance, and comparing forest versus single tree.

\medskip
\textbf{Implementation policy (strict):}
\begin{itemize}
  \item Implement the random forest core algorithm with NumPy only.
  \item Do not use a library random forest implementation for the main task.
  \item You may reuse your decision tree implementation from Lab 5.
  \item Plotting libraries are allowed for visual inspection.
  \item External libraries are allowed only for data loading if explicitly needed.
  \item Work through the notebook in order.
\end{itemize}

\section*{Core Equations And Rules}
\begin{itemize}
  \item Bootstrap sample:
  \[
    D^{(b)}
    =
    \{(x_i^{(b)},y_i^{(b)})\}_{i=1}^{n}
    \quad\text{drawn uniformly with replacement from }D.
  \]
  \item Bagging prediction (majority vote):
  \[
    \hat{y}(x)
    =
    \arg\max_{c\in\{1,\dots,C\}}
    \sum_{b=1}^{B}\mathbf{1}\{\hat{y}^{(b)}(x)=c\}.
  \]
  \item Feature subset size (classification):
  \[
    m = \lfloor\sqrt{d}\rfloor.
  \]
  \item Tree impurity and split: same as in Lab 5 (Gini impurity, impurity decrease $\Delta I$, threshold search).
  \item Out-of-bag error: each tree predicts its OOB samples; aggregate OOB predictions by majority vote per sample; compare with the true labels.
  \item Feature importance for feature $j$:
  \[
    \mathcal{I}_j
    =
    \frac{1}{B}
    \sum_{b=1}^{B}
    \sum_{\substack{m\in T_b\\ \text{split}(m)=j}}
    \frac{n_m}{n}\,\Delta I(m).
  \]
\end{itemize}

\section*{Data Files And Preparation}
\begin{itemize}
  \item \texttt{random\_forests.ipynb}: guided notebook with the assignment steps.
  \item \texttt{advanced\_ai/forest/random\_forest.py}: starter implementation.
  \item \texttt{prepare\_datasets.py}: generates the prepared data bundles.
  \item \texttt{scripts/prepare\_all\_data.sh}: shell entrypoint for dataset preparation.
  \item \texttt{scripts/check\_prepare\_all\_data.sh}: verification script for the prepared bundles.
  \item \texttt{toy\_forest.py}, \texttt{iris\_forest.py}, \texttt{oob\_and\_importance.py}: task runners.
\end{itemize}

\medskip
Expected prepared files:
\begin{itemize}
  \item \texttt{datasets/toy\_forest.npz}
  \item \texttt{datasets/iris\_forest.npz}
  \item \texttt{datasets/forest\_study.npz}
\end{itemize}

\medskip
The scripts use the conda environment \texttt{ai\_courses}.

\section*{Roadmap}
\begin{itemize}
  \item Start with Step 1: prepare and verify the datasets.
  \item After that, continue through the steps in order.
  \item Each step introduces only the task you need next.
\end{itemize}

\section*{Step 1: Prepare the Datasets}
\begin{itemize}
  \item From \texttt{assignment6\_random\_forests/}, run:
  \[
    \texttt{bash scripts/prepare\_all\_data.sh}
  \]
  \item Then run:
  \[
    \texttt{bash scripts/check\_prepare\_all\_data.sh}
  \]
  \item Confirm that all three expected \texttt{.npz} bundles exist.
\end{itemize}

\section*{Step 2: Warm-up by Hand}
\begin{itemize}
  \item Three trees predict the following votes for a test example:
  \[
    \hat{y}^{(1)} = A,\quad
    \hat{y}^{(2)} = A,\quad
    \hat{y}^{(3)} = B.
  \]
  Compute the forest prediction. Then repeat for votes $[B, B, A]$.
  \item A bootstrap sample of size $n=8$ is drawn. The sampled indices are
  \[
    [0, 1, 3, 3, 5, 6, 7, 7].
  \]
  Identify the OOB indices. Count how many distinct OOB samples there are out of $8$.
\end{itemize}

\section*{Step 3: Implement the Random Forest Core}
\begin{enumerate}[label=\textbf{T\arabic*.}]
  \item \textbf{Bootstrap sampler}
\begin{itemize}
  \item Open \texttt{advanced\_ai/forest/random\_forest.py}.
  \item Fill in \texttt{\_bootstrap\_sample}.
  \item Draw $n$ indices from $\{0,\dots,n-1\}$ with replacement.
  \item Return the array of sampled indices.
\end{itemize}

  \item \textbf{Feature subsampler}
\begin{itemize}
  \item Fill in \texttt{\_feature\_subset}.
  \item For a dataset with $d$ features, randomly select $m = \lfloor\sqrt{d}\rfloor$ feature indices.
  \item Return the array of selected feature indices.
\end{itemize}

  \item \textbf{Bagging training loop}
\begin{itemize}
  \item Fill in \texttt{fit}.
  \item For each of $B$ rounds, draw a bootstrap sample.
  \item Create a new decision tree (reuse the class from Lab 5) and train it on the bootstrapped data.
  \item Store the trained tree and the OOB index mask.
\end{itemize}

  \item \textbf{Prediction (classification)}
\begin{itemize}
  \item Fill in \texttt{predict}.
  \item Each stored tree makes a prediction for the input.
  \item Return the majority class across all trees.
\end{itemize}

  \item \textbf{Out-of-bag error}
\begin{itemize}
  \item Fill in \texttt{oob\_error}.
  \item For each training sample, find which trees had it OOB.
  \item Aggregate the votes of those trees.
  \item Compare the aggregated prediction with the true label.
  \item Return the fraction of OOB predictions that are wrong.
\end{itemize}

  \item \textbf{Feature importance}
\begin{itemize}
  \item Fill in \texttt{feature\_importances}.
  \item Walk through each tree and each internal node.
  \item Sum $\frac{n_m}{n}\Delta I$ for each feature used as the split feature.
  \item Average over $B$ trees and normalize so the importances sum to 1.
\end{itemize}
\end{enumerate}

\section*{Step 4: Run the Self-Checks}
\begin{itemize}
  \item After implementing the core class, run:
  \[
    \texttt{python -m advanced\_ai.self\_checks}
  \]
  \item Continue only after every check passes.
\end{itemize}

\section*{Step 5: Toy Dataset}
\begin{itemize}
  \item Run \texttt{toy\_forest.py}.
  \item Report the depth, number of leaves, and training accuracy of a single tree.
  \item Then report the forest accuracy (same dataset) with $B=50$.
  \item Save the decision-region comparison plot (single tree side by side with forest).
  \item Explain in one sentence why the forest boundary is smoother.
\end{itemize}

\section*{Step 6: Iris Data}
\begin{itemize}
  \item Run \texttt{iris\_forest.py}.
  \item Sweep over $B = 1, 5, 10, 20, 50, 100, 200$.
  \item Record training accuracy, OOB error, and test accuracy for each $B$.
  \item Plot OOB error and test accuracy versus $B$ on the same axes.
  \item State at which $B$ the test accuracy plateaus.
\end{itemize}

\section*{Step 7: OOB Error and Feature Importance}
\begin{itemize}
  \item Run \texttt{oob\_and\_importance.py}.
  \item Report the final OOB error and compare it with the test accuracy.
  \item List the top-3 features and their importance scores.
  \item Explain whether the top features match what you know about the dataset.
\end{itemize}

\section*{Step 8: Optional Extension}
\begin{itemize}
  \item Try different feature subset sizes: $m = 1$, $m = \lfloor\sqrt{d}/2\rfloor$, $m = d$.
  \item Report how $m$ changes the OOB error curve.
  \item Try regression with a random forest on a simple 1D dataset.
\end{itemize}

\section*{Submission}
\begin{itemize}
  \item Submit the notebook, the completed starter files, and the generated output files from the required runners.
  \item Use \texttt{collect\_submission.ipynb} or \texttt{collectSubmission.sh} to bundle the submission.
\end{itemize}
